\documentclass[11pt]{article}

\usepackage[utf8]{inputenc}
\usepackage[T1]{fontenc}
\usepackage{lmodern}
\usepackage{geometry}
\usepackage{amsmath,amssymb,amsfonts,amsthm,mathtools}
\usepackage{bm}
\usepackage{enumitem}
\usepackage{microtype}
\usepackage{hyperref}
\usepackage{titlesec}
\usepackage{booktabs}
\usepackage{array}
\usepackage{setspace}
\usepackage{mathrsfs}
\usepackage{physics}

\geometry{margin=1in}
\hypersetup{colorlinks=true, linkcolor=black, urlcolor=blue, citecolor=black}
\setlist[enumerate]{topsep=0.4em,itemsep=0.35em}
\setlist[itemize]{topsep=0.3em,itemsep=0.25em}
\titleformat{\section}{\large\bfseries}{\thesection.}{0.5em}{}
\titleformat{\subsection}{\normalsize\bfseries}{\thesubsection.}{0.5em}{}
\setstretch{1.06}

\newcommand{\Aether}{\AE{}ther}
\newcommand{\Aflow}{\AE{}ther-flow}
\newcommand{\TheoryName}{\Aflow{} Interpretation of Relativity}
\newcommand{\TheoryProgram}{\Aether{} / \Aflow{} framework}
\newcommand{\CompletedMetric}{g^{\sharp}}
\newcommand{\PreProtoSectionPackage}{\mathfrak R^{\mathrm{sub,proto,presec}}}
\newcommand{\PreProtoSectionMinSectionCore}{\mathfrak R^{\mathrm{sub,proto,presec,sec}}}
\newcommand{\PreProtoSectionResidualPackage}{\mathfrak R^{\mathrm{sub,proto,presec,res}}}
\newcommand{\ProtoSectionPackage}{\mathfrak Q^{\mathrm{sub,proto,secgen}}}

\newtheorem{definition}{Definition}
\newtheorem{proposition}{Proposition}
\newtheorem{remark}{Remark}
\newtheorem{corollary}{Corollary}
\newtheorem{theorem}{Theorem}

\title{\textbf{The \TheoryName{}}\\[0.4em]
\large Primitive-Reservoir Observer-Reduction\\
\large Proto-Section-Generating Substrate Pre-Proto-Section Dynamics\\
\large Section-Centered Normal Form Next Benchmark Criterion}
\author{}
\date{}

\begin{document}

\maketitle

\begin{abstract}
The current active-primary theorem already removed residual benchmark-facing presentation freedom from the current proto-section-generating substrate pre-proto-section dynamics package \(\PreProtoSectionPackage\) by proving that, once the fixed proto-section benchmark base is held fixed, the package is necessary and unique up to a canonical section-centered residual normal form \(\PreProtoSectionResidualPackage\). This renews the live burden again. Another same-output deeper-origin relay that merely preserves that same residual normal form and therefore merely reproduces the same current pre-proto-section package \(\PreProtoSectionPackage\), the same current observer-relevant pre-proto-section minimizer-section core \(\PreProtoSectionMinSectionCore\), and the same downstream benchmark package no longer counts as the honest default next benchmark-facing manuscript.

This note records the routing consequence of that theorem. The next honest benchmark-facing move must either derive or justify the current section-centered residual normal form \(\PreProtoSectionResidualPackage\) from still more primitive explicit substrate structure in a way that changes benchmark-facing status, or prove a genuinely smaller theorem-level collapse strictly inside that residual normal form. The benchmark boundary remains fixed throughout. The one operative metric is still exactly \(\CompletedMetric\), the ordinary matter route is unchanged, there is no second operative metric, the low-energy matter law is not enlarged, and no stronger promotion language is licensed. The positive result is therefore routing discipline at the current residual-normal-form frontier.
\end{abstract}

\section{Introduction}

The active derivational program of \TheoryName{} is no longer at the stage where the full pre-proto-section package \(\PreProtoSectionPackage\) should be treated as carrying unstructured benchmark-facing freedom. The preceding section-centered normal-form theorem already showed that, once the current proto-section benchmark base is fixed, the remaining benchmark-facing content of that package is compressed to a canonical residual normal form \(\PreProtoSectionResidualPackage\). That theorem therefore spent the looser package-level presentation as the live stopping point.

The present note records the routing consequence of that gain. It does not reopen the larger minimizer-image presentation, the already-generated minimizer-section core, or the full pre-proto-section package as if they were again independently live. It records only which kind of next manuscript would now count as an honest benchmark-facing gain once the current burden has moved to the section-centered residual normal form itself.

\paragraph{Framework and claim boundary.}
\TheoryProgram{} denotes the broader \Aether{} / \Aflow{} framework built on the claims that the \Aether{} is the underlying four-dimensional substrate of reality, the \Aflow{} is its intrinsic ordered motion, observed three-dimensional space is the local experiential slice of that deeper substrate, S-time is the experienced order of change arising from matter, light, and the \Aflow{}, and observed expansion is the three-dimensional appearance of deeper four-dimensional ordered motion. Within that framework, \TheoryName{} denotes the exact-closure theory in which the effective gravitational dynamics are adopted to be exactly Einsteinian with universal matter coupling. In the active sequence, \emph{adoption} means use of that established relativistic dynamics without claiming substrate derivation, while \emph{derivation} is reserved for a first-principles recovery from explicit substrate variables.

The present note stays strictly inside that boundary. It records only which kind of next manuscript would now count as an honest benchmark-facing gain at the current pre-proto-section residual-normal-form frontier.

\section{Fixed Inputs}

\begin{definition}[Residual-normal-form frontier inputs]
The criterion recorded here uses the following fixed inputs.
\begin{enumerate}
    \item The benchmark exact-closure package, which fixes the public one-metric GR target of the repository.
    \item The benchmark gatekeeping layer, including the recorded safeguard that blocks same-output deeper-origin relays from counting as new front-line progress once the live benchmark-relevant datum is unchanged.
    \item The earlier theorem deriving the current observer-relevant pre-proto-section minimizer-section core \(\PreProtoSectionMinSectionCore\) from the deeper current pre-proto-section package \(\PreProtoSectionPackage\).
    \item The later theorem proving that the current pre-proto-section package \(\PreProtoSectionPackage\) is itself necessary and unique up to the canonical section-centered residual normal form \(\PreProtoSectionResidualPackage\).
\end{enumerate}
\end{definition}

\begin{remark}
The present note does not replace those manuscripts. It records the routing consequence of reading them together under the active safeguard against recursive depth drift.
\end{remark}

\section{Why the Live Burden Now Sits at Residual-Normal-Form Scope}

\begin{proposition}[The live burden is renewed at residual-normal-form scope]
At the current stage of the primitive-reservoir observer-reduction line, the section-centered normal-form theorem renews the live benchmark-facing burden at the level of the canonical residual normal form \(\PreProtoSectionResidualPackage\).
\end{proposition}

\begin{proof}
That theorem changes benchmark-facing status because it removes arbitrary package-level presentation freedom from the current pre-proto-section package \(\PreProtoSectionPackage\). Once the package is shown to be necessary and unique up to the canonical residual normal form \(\PreProtoSectionResidualPackage\), the live burden no longer sits on the looser full package presentation itself. It sits on the sharper residual datum to which the package is reduced. The earlier theorem deriving the current minimizer-section core from \(\PreProtoSectionPackage\) remains preserved main-line history and routing handoff, but it no longer defines the live frontier once the current package itself has been sharpened to \(\PreProtoSectionResidualPackage\).
\end{proof}

\begin{definition}[Benchmark-neutral deeper relay at current residual-normal-form scope]
A \emph{benchmark-neutral deeper relay at current residual-normal-form scope} is a manuscript that preserves the same benchmark one-metric output, the same ordinary matter route, the same current residual normal form \(\PreProtoSectionResidualPackage\), the same current pre-proto-section package \(\PreProtoSectionPackage\), the same current observer-relevant pre-proto-section minimizer-section core \(\PreProtoSectionMinSectionCore\), the same current proto-section package \(\ProtoSectionPackage\), and the same downstream benchmark-facing consequences while merely relocating the unexplained substrate layer deeper, repackaging the same residual normal form, or re-expanding the presentation back to the full pre-proto-section package or the already-spent minimizer-image presentation without a new compression, necessity, uniqueness, or comparably sharp routing gain.
\end{definition}

\section{The Current Residual-Normal-Form Criterion}

\begin{definition}[Admissible next benchmark-facing manuscript at current residual-normal-form scope]
At the current residual-normal-form frontier, a manuscript counts as an admissible next benchmark-facing move only if it does at least one of the following:
\begin{enumerate}
    \item derives or justifies the current section-centered residual normal form \(\PreProtoSectionResidualPackage\) from still more primitive explicit substrate structure in a way that does more than merely relocate the unexplained layer deeper;
    \item proves a genuinely smaller theorem-level collapse strictly inside \(\PreProtoSectionResidualPackage\);
    \item does either of the previous two while preserving the same benchmark one-metric package, the same ordinary matter route, and the same claim boundary.
\end{enumerate}
\end{definition}

\begin{theorem}[Next honest benchmark-facing task at current residual-normal-form scope]
The next honest benchmark-facing manuscript below the current residual-normal-form frontier must either derive or justify the current section-centered residual normal form \(\PreProtoSectionResidualPackage\) from still more primitive explicit substrate structure in a benchmark-relevant way, or else prove a genuinely smaller collapse strictly inside that residual normal form. A manuscript that only repeats the same current residual datum through another deeper relay, re-expands the discussion back to the full pre-proto-section package or already-spent minimizer-image presentation without such a new gain, or invokes a quantum branch without doing the same benchmark-facing work does not satisfy the criterion.
\end{theorem}

\begin{proof}
By Proposition 1, the live burden already lies at the sharper residual normal form \(\PreProtoSectionResidualPackage\) rather than at the looser full pre-proto-section package \(\PreProtoSectionPackage\). Therefore a subsequent manuscript must improve on that status rather than merely accompany it. A deeper-origin manuscript can do so only if it changes benchmark-facing status by more than depth alone. If it simply reproduces the same current residual datum and its same downstream outputs, the routing safeguard classifies it as benchmark neutral. The remaining honest alternatives are therefore exactly those stated in the definition: either a more primitive derivation or justification of the current residual normal form itself, or a sharper internal collapse strictly inside that residual normal form.
\end{proof}

\section{What the Next Move Should Not Be}

\begin{proposition}[Screened next moves at the current residual-normal-form frontier]
At the current residual-normal-form frontier, none of the following is the default next benchmark-facing move:
\begin{enumerate}
    \item another same-output deeper-origin relay that merely preserves the current residual normal form \(\PreProtoSectionResidualPackage\) and therefore merely reproduces the same current pre-proto-section package \(\PreProtoSectionPackage\), the same current observer-relevant pre-proto-section minimizer-section core \(\PreProtoSectionMinSectionCore\), the same current proto-section package \(\ProtoSectionPackage\), and the same downstream benchmark-facing consequences;
    \item another re-expansion back to the full pre-proto-section package or to the already-spent minimizer-image presentation as if those larger presentations were again independently live burdens;
    \item a quantum branch that does not derive or justify the same residual normal form, collapse it further, or close a benchmark derivation gate in a genuinely new way;
    \item a manuscript that introduces a second operative metric, enlarges the ordinary matter law, or uses stronger promotion language.
\end{enumerate}
\end{proposition}

\begin{proof}
Item (1) fails the preceding criterion because it leaves the renewed live burden unchanged. Item (2) likewise fails because it re-expands to larger already-spent presentations rather than removing remaining freedom in \(\PreProtoSectionResidualPackage\) or proving a smaller internal datum there. Item (3) does not directly address the live burden at current scope unless it performs the same benchmark-facing work named above. Item (4) violates the fixed claim boundary of the benchmark package and therefore cannot count as a disciplined next step on the active line. The benchmark package remains the exact one-metric GR package already fixed by adoption, so any admissible next manuscript must preserve that boundary.
\end{proof}

\begin{remark}
This screening statement is not a denial that those deeper or alternative manuscripts may matter strategically. It is a routing statement: they are not the default way to spend the next benchmark-facing move from the current residual-normal-form frontier unless they do the same benchmark-facing work named above.
\end{remark}

\section{Conclusion}

The positive result of this note is routing discipline at the current residual-normal-form frontier. The section-centered normal-form theorem remains the live benchmark-facing gain because it already removes arbitrary presentational freedom from the current pre-proto-section package on the active line. The earlier theorem deriving the current minimizer-section core from that package remains preserved main-line history, but under the recursive-depth safeguard it is no longer itself the default current frontier.

The scope remains narrow and honest. The benchmark package is unchanged: the one operative metric is still exactly \(\CompletedMetric\), the ordinary matter route is unchanged, there is no second operative metric, and the low-energy matter law is not enlarged. Nothing here upgrades adoption to derivation or licenses stronger promotion language.

The next open burden is therefore clear. The next benchmark-facing manuscript should derive or justify the current section-centered residual normal form from still more primitive explicit substrate structure in a way that changes benchmark-facing status, unless the sharper honest gain available instead is a genuinely smaller theorem-level collapse strictly inside that residual normal form. Another same-output deeper relay that merely preserves that residual normal form, a re-expansion back to the larger pre-proto-section package or minimizer-image presentation, or a quantum branch that does not do that same benchmark-facing work is not the clean next manuscript target at current scope.

\end{document}
